\documentclass[10pt]{article}

\usepackage[utf8]{inputenc}

\usepackage[T1]{fontenc}

\usepackage{amsmath}

\usepackage{amsfonts}

\usepackage{amssymb}

\usepackage[version=4]{mhchem}

\usepackage{stmaryrd}

\usepackage{bbold}

\usepackage{graphicx}

\usepackage[export]{adjustbox}

\graphicspath{ {./images/} }



\begin{document}

и имеющая вид



\[

u(M, t)=\frac{\partial}{\partial t}\left\{t \Gamma_{a t}(\varphi)\right\}+t \Gamma_{a t}(\psi)

\]



где



\[

\Gamma_{a t}(\varphi)=\frac{1}{4 \pi} \int_{S_{a t}} \varphi(P) d \Omega

\]



\begin{itemize}

  \item среднее значение функции $\varphi$ на сфере $S_{a t}$ в пространстве $(x, y, z)$ радиуса at с дентром в точке $M, d \Omega-$ элемент площади единичной сферы. В случае неоднородного волнового уравнения в формуле (1) добавляется третье слагаемое (см. [2]).

\end{itemize}



Из формулы (1) спуска методом получаются формулы решения задачи Коши для случая двух (П у а с с о в а ф о р м у л а) и одного (Д'А ла м б е р a фо р м ул а) пространственного переменного (см.[2]). См. также Кирхаоба формула.



\begin{enumerate}

  \setcounter{enumi}{2}

  \item Иногда П. ф. наз. интегральное представление решения задачи Коши для уравнения теплопроводности в пространстве $\mathbb{R}^{3}$ :

\end{enumerate}



\[

\begin{gathered}

\frac{\partial u}{\partial t}-a^{2} \Delta u=0, t>0, M=(x, y, z), \\

-\infty<x, y, z<+\infty, \\

u(M, 0)=\varphi(M),

\end{gathered}

\]



имеющее вид



$u(M, t)=\frac{1}{\left(2 \sqrt{\pi a^{2} t}\right)^{3}} \int_{\mathbb{R}_{3}} \varphi(P) e^{-\frac{|M P|^{3}}{4 a^{2} t}} d \sigma(P)$.



Формула (2) непосредственно обобщается на любоө число пространственных переменных $n \geqslant 1$.



Jum.: [1] P o is s o n S. D., "Mém. Acad. sci.», 1818, t. 3,



\includegraphics[max width=\textwidth, center]{2023_07_08_e4ccd5c07d036f3d852eg-1}

д и м и р о в В. С., Уравнения математической физики, 4 изд.,



Е. Д. Соломенцев.

IIУ ССОНА ФОРМУЛА СУММИРОВАНИЯ мула



\[

\sum_{k=-\infty}^{+\infty} g(2 k \pi)=\sum_{k=-\infty}^{+\infty} \frac{1}{2 \pi} \int_{-\infty}^{+\infty} g(x) e^{-i k x} d x .

\]



П. ф. с. имеет место, если, напр., функция $g(x)$ абсолютно интегрируема на ивтервале $(-\infty,+\infty)$, имеет ограниченное изменение и $2 \mathrm{~g}(x)=g(x+0)+g(x-0)$. П. ф. с. записывается также в виде



\[

\sqrt{\bar{a}} \sum_{k=-\infty}^{+\infty} g(a k)=\sqrt{b} \sum_{k=-\infty}^{+\infty} \chi(b k),

\]



где $a$ и $b$ - любые два положительных числа, удовлетворяющие условию $a b=2 \pi$, а $\chi(u)$ есть преобразование Фурье функции $g$ :



\[

\chi(u)=\frac{1}{\sqrt{2 \pi}} \int_{-\infty}^{+\infty} g(x) e^{-i u x} d x .

\]



Лum.: [1] 3 и г м у н Д А., Тригонометршческие ряды, пер. с англ., т. 1, М., 1965; [2] Т и т ч м а p ш Е., Введение в теорию интегралов Фурье, пер. с англ., М.- Л., 1948. И. И. Волков.



ПУАССОНОВСКИЙ ПОТОК - то же, что пуассоновский процесс. Этот термин используют, как правило, в теории массового обслуживания.



ПУ АССОНОВСКИИ ПРОЦЕСС - случайный процесс $X(t)$ с независимыми прирапениями $X\left(t_{2}\right)-X\left(t_{1}\right)$, $t_{2}>t_{1}$, имеющими Пуассона распределение. В однородном П. П. Для любых $t_{2}>t_{1}$



\[

\begin{aligned}

\mathrm{P}\left\{X\left(t_{2}\right)-X\left(t_{1}\right)\right. & =k\}=\frac{\lambda^{k}\left(t_{2}-t_{1}\right)^{k}}{k !} e^{-\lambda\left(t_{2}-t_{1}\right)}, \\

k & =0,1,2, \ldots

\end{aligned}

\]



Коэффициент $\lambda>0$ наз. ин т е в с и в в о с т ь ю п уа с с оно в с к о го п р оц е с с а $X(t)$. Траектории П. п. $X(t)$ представляют собой ступенчатые функции со скачками размера 1. Моменты скачков $0<\tau_{1}<\tau_{2}<\ldots$ образуют простейщuй поток, описывающий поток тре- бований во мвогих системах массового обслуживания. Распределения случайных величин $\tau_{n}-\tau_{n-1}$ независимы при $n=1,2, \ldots$ и имеют показательную плотность $\lambda e^{-\lambda t}, \quad t \geqslant 0$.



Одним из свойств П. п. является следующее: условное распределение моментов скачков $0<\tau_{1}<\tau_{2}<$. . $<\tau_{n}<t$ при $X(t)-X(0)=n$ совпадает с распределением вариационного ряда независимой выборки объема $n$ с равномерным распределением на $[0, t]$. С другой стороны, если $0<\tau_{1}<\tau_{2}<\ldots<\tau_{n}$ - описанвый выше вариациовный ряд, то при $n \rightarrow \infty, t \rightarrow \infty$, с $n / t \rightarrow$ $\rightarrow \lambda$ получают в пределе распределение скачков П. п.



В неоднородном П. П. интенсивность $\lambda(t)$ зависит от времени $t$ и распределение $X\left(t_{2}\right)-X\left(t_{1}\right)$ определяется формулой



$\mathbf{P}\left\{X\left(t_{2}\right)-X\left(t_{1}\right)=k\right\}=\frac{\left[\int_{t_{1}}^{t_{2}} \lambda(u) d u\right]^{k}}{k !} e^{-\int t_{1}^{t_{2}} \lambda(u) d u}$.



При определенных условиях П. п. может быть показан как предел суммы неограниченно возрастающего числа независимых "редких" потоков довольно общего вида. $\mathrm{O}$ нек-рых поучительных парадоксах, связанных с П. п., см. [3], т. 2, гл. 1.



Лuт.: [ 1] Б о р о в к о в А. А., Теория вероятностей, М., $1976 ;$ [2] К., 1979; [3] $\Phi$ е л л е р В., Введение в теорию вероятностей и её приложения, 2 изд., пер. с англ., т. 1-2, М., 1967.



ПУЛЬВЕРИЗАЦИЯ н а д и ф ф е р е н и и р у ем о м м о г о о б р а и и $M$-векторное поле $W$ на касательном пространстве $T M$, имеющее в терминах локальных координат $\left(x^{1}, \ldots, x^{n}, v^{1}, \ldots, v^{n}\right)$ на $T M$, естественным образом связанных с локальными координатами $\left(x^{1}, \ldots, x^{n}\right)$ на $M$, компоненты $\left(v^{1}, \ldots, v^{n}, f^{1}, \ldots, f^{n}\right)$, где $f^{l}=f_{i}\left(x^{1}, \ldots, x^{n}, v^{1}, \ldots, v^{n}\right)-$ функции класса $C^{1}$, причем при фиксированных $x^{1}, \ldots, x^{n}$ они являются положительно однородными функциями от $v^{1}, \ldots, v^{n}$ стөпени 2 (әти свойства $W$ не зависят от конкретного выбора локальных координат). Определяемая этим полем система дифференциальных уравнений $\frac{d x^{i}}{d t}=v^{i}, \frac{d v^{i}}{d t}=f^{i}\left(x^{1}, \ldots, x^{n}, v^{1}, \ldots, v^{n}\right), i=1, \ldots, n$, эквивалентна системе дифференциальных уравнений 2-го порядка



\[

\frac{d^{2} x^{i}}{d t^{2}}=f^{i}\left(x^{1}, \ldots, x^{n}, v^{1}, \ldots, v^{n}\right),

\]



поэтому П. описывает (причем инвариантным образом, т. е. не зависящим от системы координат) систему таких уравнений на $M$.



Важнейший случай П.- когда $f^{i}$ суть многочлены 2-й степени от $v^{i}$ :



\[

f^{i}=\sum \Gamma_{j k}^{i}\left(x^{\mathbf{1}}, \ldots, x^{n}\right) v^{j_{v}}, \quad \Gamma_{j k}^{i}=\Gamma_{k j}^{i} .

\]



В әтом случае $\Gamma_{j k}^{i}$ задают на $M$ аффинную сөязность с нулевым төнзором кручения. Обратно, для всякой афффинной связности уравнения геодезич. линий задаются вөк-рой П. с fi вида $\left(^{*}\right)$ (причем при пөреходе от связности к пульверизации $\Gamma_{j k}^{i}$ симметризуются по нижним пндексам). Если поле $W-$ класса $C^{2}$, то $f^{i}$ обязавы иметь вид (*). В общем случае $W$, каким бы гладким оно ни было вне нулевого сечения расслоения $T M$, не обязано быть полем класса $C^{2}$ возле этого сечения. В такой ситуации иногда говорят об обобщенной П., оставляя термин «П.» только для специального случая $\left(^{*}\right)$. Дифференциальные уравнения для геодезических в финслеровой геометрии приводят к обобщенной $\Pi$.



Можно дать определение П. в инвариантных терминах, пригодное и для банаховых многообразий (cM. [1]).





\end{document}